\documentclass[11pt]{article}
\usepackage[margin=1in]{geometry}
\usepackage{graphicx,booktabs,amsmath,amssymb}
\usepackage{xcolor}
\usepackage[colorlinks=true,linkcolor=blue,citecolor=blue,urlcolor=blue]{hyperref}
\usepackage{authblk}

\title{\textbf{Where Does a Language Model Keep What It Knows About Relevance?\\
Attention Heads Decode Retrieval Relevance Better Than the Output Representation}}
\author{Chenhui Feng}
\affil{Nanjing University of Aeronautics and Astronautics}
\date{}

\newcommand{\auc}{\textsc{auc}}
\newcommand{\todo}[1]{\textcolor{red}{[TODO: #1]}}

\begin{document}
\maketitle

\begin{abstract}
Ask a language model whether a retrieved passage answers a query and its
zero-shot verdict is the weakest of the four relevance signals we can extract
from the same forward pass. We compare four \emph{readers} that all consume the same frozen
forward pass over the same prompt and the same answer token: the model's literal
verdict, a trained linear probe on the final-layer hidden state (the output
representation), a trained probe on the best residual-stream layer, and a probe
on the top attention heads. Across eight instruction- and base-tuned models from
three families (LLaMA, Mistral, Qwen; 3B--8B; 1500 queries each), reading the
attention heads decodes relevance better than reading the output representation
on \emph{every} model. The advantage is $+0.06$ to $+0.15$ \auc, with
query-clustered bootstrap confidence intervals that exclude zero after
Benjamini--Hochberg correction across the full $8\times4$ family of tests
($q\le0.001$ throughout)---it is the only one of the four contrasts we test
that survives correction on all eight models. The edge is not an artefact of
dimensionality: matching the final layer's read-out procedure to the head
probe's feature budget leaves it essentially unchanged. It is also not an
artefact of training: fitting a probe on the output representation never
separates from the model's own zero-shot verdict after correction (on seven of
eight models the judge is in fact the stronger of the two). We are explicit about where
the claim stops. The comparison is defined against one particular output-side
reader, and re-reading that same layer with a different read-out (a per-chunk
ensemble) closes or reverses the edge on two of eight models; and the
residual-stream version of the comparison is significant on the three base
models only. What we can say without qualification is this: for the task of
judging retrieval relevance, the attention heads of a frozen LLM are the best
place we found to read the answer, and the model's own output is the worst.
\end{abstract}

\section{Introduction}

A retrieval-augmented system hands a language model a passage and a question.
The model must decide whether the passage is relevant. In our experiments it
usually can---just not where you would look first.

Consider what the same forward pass offers. The model emits a verbal verdict,
and reading the yes/no logit difference is free but weak. A logistic regression
on the final-layer hidden state does better, because it is trained on labelled
relevance. But the final layer is only one place in a deep stack, and it is the
place furthest along a computation that must also serve every other task the
model performs. We asked a narrow, measurable question: \emph{holding the model,
the prompt, the token position, and the training budget fixed, which internal
representation is the best place to read relevance from?}

The answer, across eight models, is the attention heads. This is a
\emph{relative} result and we keep it that way. We do not claim that a probe
beats a trained cross-encoder reranker; an earlier phase of this project tested
that and found it false, and we do not revisit it. We do not claim cross-domain
generalisation. We claim something smaller and, we think, more useful: that the
choice \emph{where} to read is worth $+0.06$ to $+0.15$ \auc, reliably, across
model families and training regimes.

\paragraph{Contributions.}
\begin{itemize}
  \item \textbf{A controlled four-way comparison} (\S\ref{sec:setup}) on one
  frozen forward pass, with layer and head selection done on a held-out
  validation split, and with all four readers reporting over identical rows.
  \item \textbf{A result that replicates across the panel} (\S\ref{sec:results}):
  attention heads beat the output representation on 8/8 models, and they are the
  only one of the four contrasts that survives Benjamini--Hochberg correction
  everywhere.
  \item \textbf{Two controls that rule out the obvious explanations}
  (\S\ref{sec:controls}): a capacity-matched read-out control that leaves the
  edge intact, and a training control showing that a probe on the output
  representation adds nothing over the zero-shot verdict.
  \item \textbf{An explicit statement of where the claim stops}
  (\S\ref{sec:limits}), including the read-out sensitivity we found on two
  models and the fact that the residual-stream form of the comparison holds only
  on base models.
\end{itemize}

\paragraph{What we deliberately do not claim.}
This is a probing study, so every reader here is correlational. We initially
expected to add a causal account of \emph{why} the output is weak, and we ran
that experiment; it did not support the claim we hoped for, and we report the
null in the companion note rather than dress it up. The present paper therefore
makes no mechanistic assertion. It is about measurement.

\section{Setup}
\label{sec:setup}

\subsection{Task and data}

We use MS MARCO v1.1, forming (query, passage, label) triples. Each query
contributes its relevant passages (as marked in the dataset) and hard negatives
drawn from the same query's non-selected retrieved passages, so ``irrelevant''
means \emph{retrieved but not judged relevant} rather than \emph{random}.
Splits are made \emph{by query}, 70/15/15, so no query appears in two splits.
We sample 1500 queries, giving 7456 triples: 5218 train (1050 queries), 1121
validation (225), and 1117 test (225; 245 relevant rows).

Three disciplines keep the comparison honest. First, \textbf{layer and head
selection use the validation split only}; the test split is touched once, and
every per-layer or per-head maximum we report is a validation choice. Fields
whose names carry \texttt{DO\_NOT\_REPORT} exist to make the test-selected
alternative explicit and excluded. Second, all four readers are evaluated on
\textbf{identical rows} with identical labels; we assert both label and
query-id agreement against the activation cache rather than trusting a shared
index. Third, every reader uses the \textbf{same regularisation strength}
($C{=}1.0$), so the four columns of Table~\ref{tab:main} are mutually
consistent and their differences reconcile with the reported edges.

\subsection{Prompt and activation extraction}

Every reader sees the byte-identically identical prompt

\begin{quote}\ttfamily
Document: \{passage\}\textbackslash nQuestion: \{query\}\textbackslash nIs the document relevant to the question? Answer:
\end{quote}

\noindent and every activation is taken at the final \emph{answer token}. We
enforce this identity mechanically: a self-check script reconstructs the prompt
from each of the four code paths that build it and fails if the strings differ.

For each decoder layer $i$ we cache (a) the residual stream, i.e.\
\texttt{layers[i].output[0,-1,:]}, and (b) the per-head attention outputs,
obtained by reshaping \texttt{self\_attn.o\_proj.output[0,-1,:]} to
$[n_{\text{heads}}, d_{\text{head}}]$. All three model families in our panel
have $d_{\text{head}}=128$ (\texttt{hidden} $= n_{\text{heads}}\times128$), so
the per-head representation has the same geometry throughout. Activations are
cached in \texttt{float16}.

\subsection{Four readers}

\begin{enumerate}
  \item \textbf{LLM-judge} (\emph{zero training}). The normalised
  $P(\text{yes})$ vs.\ $P(\text{no})$ at the answer token, summed over five
  surface variants of each word (\texttt{yes}/\texttt{Yes}/\texttt{YES}/%
  \texttt{" yes"}/\texttt{" YES"}). No labels are used.
  \item \textbf{Output-representation probe}. An $\ell_2$-regularised logistic
  regression on the final-layer residual, fitted on train.
  \item \textbf{Internal-residual probe}. The same probe on the
  validation-selected residual layer $L^\ast$.
  \item \textbf{Attention-head probe} (\emph{ours}). Per-head validation \auc{}
  ranks all $n_{\text{layers}}\times n_{\text{heads}}$ heads; we take the top
  $k{=}20$, standardise each head's activations, concatenate, and fit one
  logistic regression on train.
\end{enumerate}

The head probe is deliberately simple. We are not proposing an architecture; we
are asking which of four equally simple ways of reading the same forward pass
recovers the most relevance. A head-space budget of $20\times128=2560$
dimensions is \emph{smaller} than the residual stream (3072--4096 dimensions),
so the head probe is not winning on feature count (Table~\ref{tab:geom}).

\begin{table}[t]
\centering\small
\caption{Feature geometry. The head probe uses fewer features than the residual
probe on every model, so its advantage cannot come from dimensionality.}
\label{tab:geom}
\begin{tabular}{lccccc}
\toprule
Model family & layers & heads & $d_{\text{head}}$ & hidden & head probe ($k{=}20$) \\
\midrule
LLaMA-3.2-3B        & 28 & 24 & 128 & 3072 & 2560 \\
LLaMA-3.1-8B        & 32 & 32 & 128 & 4096 & 2560 \\
Mistral-7B          & 32 & 32 & 128 & 4096 & 2560 \\
Qwen2.5-7B          & 28 & 28 & 128 & 3584 & 2560 \\
\bottomrule
\end{tabular}
\end{table}

\subsection{Statistics}

Because the split is by query, the resampling unit is the query. We report 95\%
confidence intervals and one-sided $p$-values from a \emph{query-clustered
paired bootstrap} (5000 resamples): the same resampled query set is used for
both readers in a contrast, which is what makes the interval a statement about
their difference rather than about two independent estimates. With 5000
resamples the smallest attainable $p$ is $2\times10^{-4}$, and we write
$p<2\times10^{-4}$ rather than $0$.

We test four contrasts over eight models, so we report Benjamini--Hochberg
$q$-values computed within each contrast family (8 tests). The headline claim of
this paper is the one that survives that correction everywhere.

Like every result here, the statistics are derived by scripts from raw JSON
artifacts, and a single collection script rebuilds the paper's number table with
a provenance path recorded for every value. We mention this because our own
earlier draft of this work quoted two different values for the same quantity in
the same table; the remedy was mechanical rather than editorial.

\section{Results}
\label{sec:results}

\subsection{The attention heads are the best reader, on every model}

Table~\ref{tab:main} and Figure~\ref{fig:readers} give the full picture. Two
things stand out.

First, \textbf{the zero-shot verdict is the weakest reader on all eight
models}. It ranges from 0.545 (LLaMA-3.2-3B) to 0.660 (Qwen2.5-7B-Instruct) and
is below the head probe everywhere. It is \emph{not}, however, below the trained
output-representation probe: the judge is the stronger of the two on seven of
eight models and the weaker on one (Mistral-7B, 0.579 vs.\ 0.583). Training a
probe on the final layer buys very little over asking the model---which is the
point of the \texttt{final $-$ judge} row in Table~\ref{tab:edges}---and that is
precisely why the interesting contrast is the \emph{head} probe against the
output representation, not the judge against anything. A model that can decode
relevance at 0.67--0.73 from its attention heads is answering at 0.55--0.66 when
asked directly.

Second, \textbf{the attention-head probe is the best reader on all eight
models}. It beats the output-representation probe by $+0.059$ to $+0.148$ \auc,
and it beats the best residual-stream layer on all eight as well, though that
second margin is smaller.

\begin{table}[t]
\centering\small
\caption{Four readers on the same frozen forward pass (test \auc, 1500 queries
per model; layer and head selection on validation; one regularisation strength
$C{=}1.0$ for all readers). \textbf{Bold}: best reader in the row. The
attention-head probe is best on all eight models. Figure~\ref{fig:readers}
shows the same numbers with confidence intervals.}
\label{tab:main}
\begin{tabular}{lcccccc}
\toprule
Model & $L$ & judge & final & best-resid & attn (ours) \\
\midrule
LLaMA-3.2-3B (base)       & 28 & 0.545 & 0.536 & 0.613 (L9)  & \textbf{0.684} \\
LLaMA-3.1-8B (base)       & 32 & 0.574 & 0.562 & 0.621 (L8)  & \textbf{0.677} \\
Mistral-7B-v0.3 (base)    & 32 & 0.579 & 0.583 & 0.606 (L16) & \textbf{0.684} \\
LLaMA-3.2-3B-Instruct     & 28 & 0.614 & 0.612 & 0.638 (L18) & \textbf{0.712} \\
LLaMA-3.1-8B-Instruct     & 32 & 0.630 & 0.596 & 0.609 (L12) & \textbf{0.725} \\
Mistral-7B-Instruct-v0.3  & 32 & 0.646 & 0.596 & 0.634 (L13) & \textbf{0.713} \\
Qwen2.5-7B (base)         & 28 & 0.652 & 0.592 & 0.655 (L22) & \textbf{0.680} \\
Qwen2.5-7B-Instruct       & 28 & 0.660 & 0.610 & 0.630 (L23) & \textbf{0.669} \\
\bottomrule
\end{tabular}
\end{table}

\begin{figure}[t]
\centering
\includegraphics[width=\textwidth]{figures/fig1_readers.png}
\caption{The four readers across eight models, with query-clustered 95\%
confidence intervals. The head probe (red) is highest in every group; the
zero-shot judge (light grey) is lowest or near-lowest everywhere.}
\label{fig:readers}
\end{figure}

\subsection{The edge survives correction, and it is the only one that does}

Figure~\ref{fig:forest} shows the head-probe-minus-output-representation edge
per model, with intervals and corrected $q$-values. It lies between $+0.059$ and
$+0.148$ and excludes zero after correction on all eight models.

Table~\ref{tab:edges} puts that in context by reporting all four contrasts we
tested. Only one survives on every model. The residual-stream contrast
(best-resid $-$ final) survives on the three base models and loses significance
on their instruct counterparts; the contrast against the judge
(best-resid $-$ judge) excludes zero on two models and survives correction on
one; and the ``pure training''
contrast (final $-$ judge) survives on none. The last of these is worth
emphasising, because it is the control that matters most for the framing of this
paper: \emph{training a probe on the output representation does not beat simply
asking the model}, on any model in our panel.

\begin{table}[t]
\centering\small
\caption{All four contrasts, with Benjamini--Hochberg $q$ within each family
(8 tests). A cell is marked $\bullet$ when the 95\% CI excludes zero after
correction. The attention-head contrast is the only one that survives
everywhere.}
\label{tab:edges}
\begin{tabular}{lcccc}
\toprule
Contrast & survives (of 8) & range of effect & models where it survives \\
\midrule
\textbf{attn $-$ final}      & \textbf{8} & $+0.059 \ldots +0.148$ & all eight \\
best-resid $-$ final         & 3 & $+0.059 \ldots +0.078$ & LLaMA-3.2-3B, LLaMA-3.1-8B, Qwen2.5-7B (base) \\
best-resid $-$ judge         & 2 (1 after BH) & $+0.047 \ldots +0.068$ & LLaMA-3.2-3B, LLaMA-3.1-8B \\
final $-$ judge              & 0 & $-0.060 \ldots +0.004$ & --- \\
\bottomrule
\end{tabular}
\end{table}

\begin{figure}[t]
\centering
\includegraphics[width=0.86\textwidth]{figures/fig2_forest.png}
\caption{The attention-head edge per model. Filled markers are base models,
hollow are instruct. Every interval clears zero after correction; the smallest
effect (Qwen2.5-7B-Instruct, $+0.059$) is also the one with the largest
correction penalty.}
\label{fig:forest}
\end{figure}

Figure~\ref{fig:pair} presents the same comparison two other ways: as a paired
slope chart, and as a scatter in which all eight models sit above the diagonal.
The pairing matters---the two readers see the same forward pass on the same rows,
so the right question is whether the difference is positive, and it is, on every
model.

\begin{figure}[t]
\centering
\includegraphics[width=\textwidth]{figures/fig4_attn_vs_final.png}
\caption{(a) Paired slope from the output representation to the attention heads;
every line rises. (b) The same eight points as a scatter; all sit in the shaded
region where the head probe wins. Colours are model families.}
\label{fig:pair}
\end{figure}

\subsection{Why is the output reader not simply under-trained?}

Two objections could explain the head advantage without any claim about
\emph{where} relevance lives, and we tested both.

\paragraph{Objection 1: the head probe gets a better read-out.}
The head probe selects 20 units on validation and concatenates them; the
residual probe reads one 3072--4096-dimensional vector. The natural control is to
give the final layer the \emph{identical} procedure: same 2560-dimensional
budget, validation-selected 20 chunks, per-unit standardisation, concatenation.
Under that matched procedure the edge barely moves ($+0.043$ to $+0.128$,
significant on 8/8). The head advantage is therefore not bought by the read-out
procedure or by the selection step---but this control cuts both ways, and we
report the other direction honestly in \S\ref{sec:limits}.

\paragraph{Objection 2: the head probe is trained and the judge is not.}
This is correct as stated, and it is why we also compare against a
\emph{trained} reader on the output representation. That reader is the
\texttt{final} column of Table~\ref{tab:main}: trained on 5218 labelled pairs,
it still loses to the heads on every model. The ``final $-$ judge'' row of
Table~\ref{tab:edges} quantifies the training contribution on its own: after
correction it is not distinguishable from zero anywhere, and its point estimates
are negative on seven of eight models. Training on the output representation, in
other words, is not what produces the advantage---the location is.

\subsection{What the depth profile looks like}

Figure~\ref{fig:depth} shows the per-layer picture for LLaMA-3.2-3B, and the
same figure for all eight models is in the artifact bundle
(\texttt{fig3\_depth\_<model>.png}). Three observations.

Relevance decodability from the residual stream rises into the middle of the
network and is no longer rising at the output layer. The peak is a
validation-selected layer, and it is well before the end (L9 of 28 for
LLaMA-3.2-3B; the peak layers for the other models are in Table~\ref{tab:main}).

The per-head map is sparse and mid-depth. The best single head reaches
$0.64$--$0.65$ on LLaMA-3.2-3B and Mistral-7B, while the \emph{mean} head in the
same layer sits at $0.50$--$0.62$: most heads carry little and a few carry a
lot, which is why a probe that selects units on validation does well here, and
why we report the selection procedure explicitly. Early layers are markedly
weaker: the first layer is exactly uninformative for every single head on
LLaMA-3.2-3B, and the first four are on Mistral-7B, whereas Qwen2.5-7B has
informative heads from the start.

The selected heads cluster at a characteristic relative depth per family
(Figure~\ref{fig:layers}): mid-network for LLaMA ($\approx$0.35--0.52 of
depth), later for Qwen (whose best single head reaches $0.72$ at layer 18 of
28), and Mistral's informative heads appear after its uninformative early band
even though its residual stream carries the signal from much earlier. We report
this as description; we did not find a principled account of why the depths
differ.

\begin{figure}[t]
\centering
\includegraphics[width=\textwidth]{figures/fig3_depth_llama32_3b.png}
\caption{LLaMA-3.2-3B. (a) Residual-stream decodability by depth, against the
zero-shot judge baseline; the peak is validation-selected, the output layer is
marked. (b) Per-head decodability; the informative heads form a sparse,
mid-depth band. (c) Best and mean head per layer versus the residual stream at
the same layer. Per-model versions are in the artifact bundle.}
\label{fig:depth}
\end{figure}

\begin{figure}[t]
\centering
\includegraphics[width=0.9\textwidth]{figures/fig5_layer_dist.png}
\caption{Where the validation-selected heads live, in relative depth. Marker
area is the number of selected heads in that layer.}
\label{fig:layers}
\end{figure}

\section{Controls and robustness}
\label{sec:controls}

\paragraph{Read-out matched to the head probe.} Described in \S\ref{sec:results};
the edge persists at $+0.043$ to $+0.128$ on 8/8 models when the final layer gets
the head probe's exact procedure.

\paragraph{Read-out replaced by an ensemble.} We also replaced the single linear
read-out of the final layer with a per-chunk logistic-regression ensemble. This
\emph{raises} the final layer substantially (e.g.\ $0.536\to0.621$ on
LLaMA-3.2-3B, a gain of $+0.085$---comparable to the entire residual-stream
edge). On two of eight models it closes or reverses the head advantage
(Qwen2.5-7B: $0.683$ vs.\ $0.680$; Qwen2.5-7B-Instruct: $0.719$ vs.\ $0.669$).
We therefore do not claim that \emph{attention heads beat the output end}; we
claim the specific matched comparison of Table~\ref{tab:main}, and we flag the
read-out sensitivity as a boundary in \S\ref{sec:limits}.

\paragraph{MLP read-outs.} At matched parameter count and at a wider setting,
final-layer MLPs reach $0.594$ and $0.560$ on LLaMA-3.2-3B---above the linear
final-layer probe ($0.536$) but still below the head probe ($0.684$).

\paragraph{Prompt wording.} The comparison holds under two rewritings of the
prompt (a terse ``Passage/Query/Relevant?'' form and a verbose form). We report
this as a direction check rather than as a strength claim: the verbose form is
$\approx$17\% longer in tokens, and we did not attempt to separate wording from
length.

\paragraph{A second dataset.} On FiQA with BM25 hard negatives (LLaMA-3.2-3B,
1500 queries), the head probe reaches $0.884$ against a $0.724$ output
representation and a $0.597$ judge, so the ordering is unchanged on a second
in-domain dataset. This is one model, not eight, and we do not lean on it.

\section{Where the claim stops}
\label{sec:limits}

\paragraph{This is one specific comparison.} The headline contrast is defined
against a single linear read-out of the final-layer residual. We showed in
\S\ref{sec:controls} that a different read-out of that same layer changes the
answer on two of eight models. A reader who prefers a per-chunk ensemble as
``the'' output representation should read our result as
\emph{read-out-dependent}, not as a general claim about the output end.

\paragraph{The residual-stream form of the claim is narrower.} ``Internal
beats output representation'' is often stated for the residual stream. In our
panel that contrast is significant on the three base models only, and its
instruct counterparts show smaller, non-significant effects. We report the head
comparison as the headline precisely because it is the one that replicates
everywhere.

\paragraph{The judge is a zero-shot judge.} We compare against the model's own
verdict with no exemplars. A stronger prompt or a few-shot judge would raise the
lower bound. We tested one such variant in an earlier phase and it did not close
the gap, but we do not have that artifact in the current bundle, so we state the
comparison as ``zero-shot verdict'' and nothing more.

\paragraph{Correlational, not causal.} Every reader here is a probe. The
companion work on this project attempted a causal account of the output's
weakness and did not support it; the present paper therefore contains no
mechanistic assertion, and readers should not infer one.

\paragraph{Relevance means MS MARCO relevance.} Labels come from dataset qrels
for one retrieval collection (plus one FiQA replication). We make no
cross-domain and no reranking-accuracy claim.

\section{Related work}
\label{sec:related}

\todo{restore citations: this section is deliberately citation-free in this
draft. The bibliography in \texttt{paper/refs.bib} already contains the entries
this project verified as real, covering (i) probing classifiers,
(ii) the logit/tuned lens, (iii) activation steering and representation
engineering, (iv) calibration and hallucination work showing models encode
correctness they do not surface, and (v) concurrent work on hallucination
detection from mid-layer hidden states. The net-novelty statement to write is:
internals-outstrip-the-output is not new as a slogan, but a controlled,
corrected, eight-model comparison localised to \emph{retrieval relevance}, with
the head-versus-output-representation contrast as the surviving effect, is.}

\section{Conclusion}

Across eight models from three families, reading the top attention heads of a
frozen forward pass recovers retrieval relevance better than reading the
final-layer hidden state, better than any single residual layer, and far better
than the model's own verdict. The margin over the output representation is
$+0.06$ to $+0.15$ \auc, it holds on every model, and it is the only one of four
contrasts we tested that survives multiple-comparison correction everywhere.
Training is not the explanation---a probe on the output representation never
beats simply asking the model. If you need to know whether a passage is
relevant, the model's heads know, and its mouth is the last place to look.

\appendix
\section{Artifact map}

Every number in this paper traces to a JSON artifact, and every figure is
regenerated by a script from those artifacts.

\begin{table}[h]
\centering\small
\begin{tabular}{ll}
\toprule
Artifact & Contents \\
\midrule
\texttt{results/PAPER\_NUMBERS.json} & single source of truth, with provenance per value \\
\texttt{results/significance\_<model>\_q1500.json} & four reader AUCs + CIs + the four edges \\
\texttt{results/significance\_summary\_q1500.json} & BH-corrected edge table across models \\
\texttt{results/capacity\_control\_<model>\_q1500.json} & read-out matched / ensemble / MLP controls \\
\texttt{results/figures/fig3\_depth\_<model>.png} & per-model depth + per-head figure \\
\texttt{results/figures/fig*.json} & provenance sidecar for every figure \\
\bottomrule
\end{tabular}
\end{table}

\noindent Scripts: \texttt{scripts/collect\_paper\_numbers.py} (paper numbers),
\texttt{scripts/aggregate\_significance.py} (BH correction),
\texttt{scripts/plot\_paper\_figures.py} (all figures),
\texttt{scripts/selfcheck.py} (structural self-check, including that all four
prompt builders agree byte-for-byte).

\end{document}
